%%%%%%%%%%%%%%%%%%%%%%%%%%%%%%%%%%%%%%%%%
% FRI Data Science_report LaTeX Template
% Version 1.0 (28/1/2020)
% 
% Jure Demšar (jure.demsar@fri.uni-lj.si)
%
% Based on MicromouseSymp article template by:
% Mathias Legrand (legrand.mathias@gmail.com) 
% With extensive modifications by:
% Antonio Valente (antonio.luis.valente@gmail.com)
%
% License:
% CC BY-NC-SA 3.0 (http://creativecommons.org/licenses/by-nc-sa/3.0/)
%
%%%%%%%%%%%%%%%%%%%%%%%%%%%%%%%%%%%%%%%%%


%----------------------------------------------------------------------------------------
%	PACKAGES AND OTHER DOCUMENT CONFIGURATIONS
%----------------------------------------------------------------------------------------
\documentclass[fleqn,moreauthors,10pt]{ds_report}
\usepackage[english]{babel}

\graphicspath{{fig/}}




%----------------------------------------------------------------------------------------
%	ARTICLE INFORMATION
%----------------------------------------------------------------------------------------

% Header
\JournalInfo{FRI Natural language processing course 2026}

% Interim or final report
\Archive{Project report} 
%\Archive{Final report} 

% Article title
\PaperTitle{Natural language processing course} 

% Authors (student competitors) and their info
\Authors{Angel Naumov, Naum Spasev, Marija Koshtrevska}

% Advisors
\affiliation{\textit{Advisors: Slavko Žitnik}}

% Keywords
\Keywords{Domain-Specific AI, Natural Language Processing, Retrieval-Augmented Generation, Game Analytics}
\newcommand{\keywordname}{Keywords}


%----------------------------------------------------------------------------------------
%	ABSTRACT
%----------------------------------------------------------------------------------------

\Abstract{
Large language models have difficulty providing accurate, up-to-date answers in rapidly evolving fields like competitive gaming. This project introduces a domain-specific AI assistant for League of Legends that uses Retrieval-Augmented Generation (RAG). The system indexes patch notes, champion statistics, and match history data and uses this information to retrieve relevant context at query time and generate informed, accurate responses. The assistant can answer questions about champion strengths and weaknesses and predict the viability of champion picks based on retrieved patch and matchup data. Evaluation combines automated retrieval metrics with human assessments of response quality and usefulness. The expected results will demonstrate that a domain-specific RAG reduces inaccurate responses and improves precision compared to a standalone LLM.


}

%----------------------------------------------------------------------------------------

\begin{document}

% Makes all text pages the same height
\flushbottom 

% Print the title and abstract box
\maketitle 

% Removes page numbering from the first page
\thispagestyle{empty} 

%----------------------------------------------------------------------------------------
%	ARTICLE CONTENTS
%----------------------------------------------------------------------------------------

\section*{Introduction}


	In recent years, large language models (LLMs) have demonstrated impressive abilities in understanding and generating human-like text across various fields. However, despite their versatility, these models struggle with tasks requiring the latest specialized knowledge or a deep understanding of context. 

These limitations are noticeable in competitive gaming environments, such as League of Legends. This complex multiplayer game has a continuously shifting strategic landscape due to frequent patch updates. The choices players make are based on a complex interaction of various factors, including the strengths and weaknesses of the champion, the items chosen, and past performance trends. This is something that a model based on statistics struggles to keep up with.

This work addresses these challenges by developing a domain-specific conversational assistant based on a RAG framework. Rather than relying on knowledge encoded during pretraining, the system dynamically retrieves relevant information from a tailored knowledge base containing patch notes, champion statistics, and match history data. This context is then incorporated directly into response generation. This design is particularly well-suited to League of Legends: when a new patch is released, only the knowledge base requires updating, not the model itself. The result is a system that can accurately answer nuanced questions about champion performance, recent balance changes, and matchup recommendations, surpassing the capabilities of standalone LLMs.

\section*{Related Work}

Earlier machine learning studies in League of Legends have mostly concentrated on match outcome prediction and player performance assessment, attaining robust outcomes through gradient boosting, tensor factorization, and event-driven win probability models~\cite{ferreira2025quantifying, maymin2021smart, costa2021feature}. Jiang et al.~\cite{jiang2020individualized} notably demonstrated that incorporating patch version context directly into predictive models improves accuracy. This confirms that game knowledge is inherently version-dependent and rapidly evolving.

However, these models primarily operate on structured data and lack mechanisms for interactive knowledge access. This has led to growing interest in NLP approaches.

Lee et al.~\cite{lee2026leaguebot} introduced a voice-based large language model companion for beginner League of Legends players in their 2026 study, LeagueBot, reporting significant reductions in cognitive challenge and perceived tension in a controlled experiment with 33 players. However, participants reported receiving recommendations for nonexistent items. The authors attributed this hallucination problem to the LLM's lack of access to current game data. RAG was explicitly proposed as a solution, but latency concerns in fast-paced gameplay prevented its implementation. Motivated by these limitations, this work presents an RAG-based assistant for League of Legends that retrieves current patch notes, champion statistics, and match data, making it possible to provide practical, up-to-date responses.
%------------------------------------------------

% \section*{Methods}

\section*{Data}

We leverage the official Riot API\cite{riotApi} alongside player data sourced from OP.GG\cite{opgg} as our data. Beyond in-game performance metrics, we integrate official documentation and community knowledge to guide feature engineering and model training. Additionally, we utilize balance patch releases from Riot Games to capture and account for evolving gameplay mechanics across different patch versions.

\subsection*{Game Rules and Domain Knowledge}

The foundational layer of the knowledge base consists of official and community-maintained documentation describing the rules, objectives, and mechanics of League of Legends. These sources provide structured, stable textual descriptions of core concepts such as champion roles, map objectives, itemisation, and win conditions. Since this content changes infrequently, it serves as a static background knowledge layer, ensuring consistent understanding of game fundamentals.

\subsection*{Beginner and Strategy Guides}

To ensure the assistant can respond to practical, player oriented queries in natural language, beginner and strategy guides are incorporated. These documents reflect how real players formulate questions and what useful answers look like in practice. This category also informs the construction of the evaluation question set, as the phrasing and topics closely mirror the queries the assistant is expected to handle.

\subsection*{Patch Notes}

Patch notes represent the primary and most critical component of the RAG knowledge base. Riot Games releases balance updates approximately every two weeks, introducing champion buffs, nerfs, item changes, and systemic adjustments that fundamentally alter the competitive landscape. Each patch document is chunked, embedded, and indexed into the vector store, allowing the retrieval pipeline to surface the most relevant balance changes in response to user queries. This design enables the knowledge base to be updated with each new patch without any modification to the underlying model.

\subsection*{Champion Statistics and Meta Data}

Champion performance statistics are computed from match data collected via the Riot API, calculating win rates, pick and ban rates, optimal item builds on a per-patch and per-rank basis. This data is inherently dynamic. A champion considered strong in one patch may be significantly weaker following a targeted nerf, making it well-suited for RAG retrieval. When a user asks whether a specific champion is viable in the current meta, the retrieval pipeline surfaces the most recent computed statistics, grounding responses in empirical data derived from observed gameplay.

\subsection*{Match Data}

Historical match data is obtained via the Riot Games Developer API, which provide granular per-game information including champion selections, item builds, kill events, gold progression, and final outcomes. This data serves primarily as ground truth for evaluation: champion pick recommendations can be assessed against actual match outcomes from Diamond+ tier games to quantify predictive accuracy.

% Use the Methods section to describe what you did an how you did it -- in what way did you prepare the data, what algorithms did you use, how did you test various solutions ... Provide all the required details for a reproduction of your work.

% Below are \LaTeX examples of some common elements that you will probably need when writing your report (e.g. figures, equations, lists, code examples ...).


% \subsection*{Equations}

% You can write equations inline, e.g. $\cos\pi=-1$, $E = m \cdot c^2$ and $\alpha$, or you can include them as separate objects. The Bayes’s rule is stated mathematically as:

% \begin{equation}
% 	P(A|B) = \frac{P(B|A)P(A)}{P(B)},
% 	\label{eq:bayes}
% \end{equation}

% where $A$ and $B$ are some events. You can also reference it -- the equation \ref{eq:bayes} describes the Bayes's rule.

% \subsection*{Lists}

% We can insert numbered and bullet lists:

% % the [noitemsep] option makes the list more compact
% \begin{enumerate}[noitemsep] 
% 	\item First item in the list.
% 	\item Second item in the list.
% 	\item Third item in the list.
% \end{enumerate}

% \begin{itemize}[noitemsep] 
% 	\item First item in the list.
% 	\item Second item in the list.
% 	\item Third item in the list.
% \end{itemize}

% We can use the description environment to define or describe key terms and phrases.

% \begin{description}
% 	\item[Word] What is a word?.
% 	\item[Concept] What is a concept?
% 	\item[Idea] What is an idea?
% \end{description}


% \subsection*{Random text}

% This text is inserted only to make this template look more like a proper report. Lorem ipsum dolor sit amet, consectetur adipiscing elit. Etiam blandit dictum facilisis. Lorem ipsum dolor sit amet, consectetur adipiscing elit. Interdum et malesuada fames ac ante ipsum primis in faucibus. Etiam convallis tellus velit, quis ornare ipsum aliquam id. Maecenas tempus mauris sit amet libero elementum eleifend. Nulla nunc orci, consectetur non consequat ac, consequat non nisl. Aenean vitae dui nec ex fringilla malesuada. Proin elit libero, faucibus eget neque quis, condimentum laoreet urna. Etiam at nunc quis felis pulvinar dignissim. Phasellus turpis turpis, vestibulum eget imperdiet in, molestie eget neque. Curabitur quis ante sed nunc varius dictum non quis nisl. Donec nec lobortis velit. Ut cursus, libero efficitur dictum imperdiet, odio mi fermentum dui, id vulputate metus velit sit amet risus. Nulla vel volutpat elit. Mauris ex erat, pulvinar ac accumsan sit amet, ultrices sit amet turpis.

% Phasellus in ligula nunc. Vivamus sem lorem, malesuada sed pretium quis, varius convallis lectus. Quisque in risus nec lectus lobortis gravida non a sem. Quisque et vestibulum sem, vel mollis dolor. Nullam ante ex, scelerisque ac efficitur vel, rhoncus quis lectus. Pellentesque scelerisque efficitur purus in faucibus. Maecenas vestibulum vulputate nisl sed vestibulum. Nullam varius turpis in hendrerit posuere.


% \subsection*{Figures}

% You can insert figures that span over the whole page, or over just a single column. The first one, \figurename~\ref{fig:column}, is an example of a figure that spans only across one of the two columns in the report.

% \begin{figure}[ht]\centering
% 	\includegraphics[width=\linewidth]{single_column.pdf}
% 	\caption{\textbf{A random visualization.} This is an example of a figure that spans only across one of the two columns.}
% 	\label{fig:column}
% \end{figure}

% On the other hand, \figurename~\ref{fig:whole} is an example of a figure that spans across the whole page (across both columns) of the report.

% % \begin{figure*} makes the figure take up the entire width of the page
% \begin{figure*}[ht]\centering 
% 	\includegraphics[width=\linewidth]{whole_page.pdf}
% 	\caption{\textbf{Visualization of a Bayesian hierarchical model.} This is an example of a figure that spans the whole width of the report.}
% 	\label{fig:whole}
% \end{figure*}


% \subsection*{Tables}

% Use the table environment to insert tables.

% \begin{table}[hbt]
% 	\caption{Table of grades.}
% 	\centering
% 	\begin{tabular}{l l | r}
% 		\toprule
% 		\multicolumn{2}{c}{Name} \\
% 		\cmidrule(r){1-2}
% 		First name & Last Name & Grade \\
% 		\midrule
% 		John & Doe & $7.5$ \\
% 		Jane & Doe & $10$ \\
% 		Mike & Smith & $8$ \\
% 		\bottomrule
% 	\end{tabular}
% 	\label{tab:label}
% \end{table}


% \subsection*{Code examples}

% You can also insert short code examples. You can specify them manually, or insert a whole file with code. Please avoid inserting long code snippets, advisors will have access to your repositories and can take a look at your code there. If necessary, you can use this technique to insert code (or pseudo code) of short algorithms that are crucial for the understanding of the manuscript.

% \lstset{language=Python}
% \lstset{caption={Insert code directly from a file.}}
% \lstset{label={lst:code_file}}
% \lstinputlisting[language=Python]{code/example.py}

% \lstset{language=R}
% \lstset{caption={Write the code you want to insert.}}
% \lstset{label={lst:code_direct}}
% \begin{lstlisting}
% import(dplyr)
% import(ggplot)

% ggplot(diamonds,
% 	   aes(x=carat, y=price, color=cut)) +
%   geom_point() +
%   geom_smooth()
% \end{lstlisting}

% %------------------------------------------------

% \section*{Results}

% Use the results section to present the final results of your work. Present the results in a objective and scientific fashion. Use visualisations to convey your results in a clear and efficient manner. When comparing results between various techniques use appropriate statistical methodology.

% \subsection*{More random text}

% This text is inserted only to make this template look more like a proper report. Lorem ipsum dolor sit amet, consectetur adipiscing elit. Etiam blandit dictum facilisis. Lorem ipsum dolor sit amet, consectetur adipiscing elit. Interdum et malesuada fames ac ante ipsum primis in faucibus. Etiam convallis tellus velit, quis ornare ipsum aliquam id. Maecenas tempus mauris sit amet libero elementum eleifend. Nulla nunc orci, consectetur non consequat ac, consequat non nisl. Aenean vitae dui nec ex fringilla malesuada. Proin elit libero, faucibus eget neque quis, condimentum laoreet urna. Etiam at nunc quis felis pulvinar dignissim. Phasellus turpis turpis, vestibulum eget imperdiet in, molestie eget neque. Curabitur quis ante sed nunc varius dictum non quis nisl. Donec nec lobortis velit. Ut cursus, libero efficitur dictum imperdiet, odio mi fermentum dui, id vulputate metus velit sit amet risus. Nulla vel volutpat elit. Mauris ex erat, pulvinar ac accumsan sit amet, ultrices sit amet turpis.

% Phasellus in ligula nunc. Vivamus sem lorem, malesuada sed pretium quis, varius convallis lectus. Quisque in risus nec lectus lobortis gravida non a sem. Quisque et vestibulum sem, vel mollis dolor. Nullam ante ex, scelerisque ac efficitur vel, rhoncus quis lectus. Pellentesque scelerisque efficitur purus in faucibus. Maecenas vestibulum vulputate nisl sed vestibulum. Nullam varius turpis in hendrerit posuere.

% Nulla rhoncus tortor eget ipsum commodo lacinia sit amet eu urna. Cras maximus leo mauris, ac congue eros sollicitudin ac. Integer vel erat varius, scelerisque orci eu, tristique purus. Proin id leo quis ante pharetra suscipit et non magna. Morbi in volutpat erat. Vivamus sit amet libero eu lacus pulvinar pharetra sed at felis. Vivamus non nibh a orci viverra rhoncus sit amet ullamcorper sem. Ut nec tempor dui. Aliquam convallis vitae nisi ac volutpat. Nam accumsan, erat eget faucibus commodo, ligula dui cursus nisi, at laoreet odio augue id eros. Curabitur quis tellus eget nunc ornare auctor.


% %------------------------------------------------

% \section*{Discussion}

% Use the Discussion section to objectively evaluate your work, do not just put praise on everything you did, be critical and exposes flaws and weaknesses of your solution. You can also explain what you would do differently if you would be able to start again and what upgrades could be done on the project in the future.


% %------------------------------------------------

% \section*{Acknowledgments}

% Here you can thank other persons (advisors, colleagues ...) that contributed to the successful completion of your project.


%----------------------------------------------------------------------------------------
%	REFERENCE LIST
%----------------------------------------------------------------------------------------
\bibliographystyle{unsrt}
\bibliography{report}


\end{document}